%% Entorno Descripci�n. En tres sabores:
    %% 1. El m�s general
    \newenvironment{descripcion}[2]%
    {\begin{basedescript}{\desclabelwidth{#2}}%
    \item[#1]}%
    {\end{basedescript}}%

    %% 2. Como poner en negrita (tambi�n parecido a section* pero puedo seguir en la misma l�nea)
    \newenvironment{descripcion0}[1]%
    {\begin{basedescript}{\desclabelwidth{-1ex}}%
    \item[#1]}%
    {\end{basedescript}}%

    %% 3. Con 1.25cm de tama�o
    \newenvironment{descripcion1}[1]%
    {\begin{basedescript}{\desclabelwidth{1.25cm}}%
    \item[#1]}%
    {\end{basedescript}}%

    %% Entorno para citar texto, 80% del tama�o (10% de cada lado)
      \newenvironment{citetext}{%
      \def\leftmargini{0.1\textwidth}%
      \def\rightmargini{0.1\textwidth}%
      \vspace*{-0.35cm}%
      \begin{quotation}\parskip 0.15cm\guillemotleft}%
      {\guillemotright\end{quotation}\vspace*{0.1cm}}


      %% Sin � (caso especial de citetext)
      \newenvironment{citetext_sin_right}{%
      \def\leftmargini{0.1\textwidth}%
      \def\rightmargini{0.1\textwidth}%
      \vspace*{-0.35cm}%
      \begin{quotation}\parskip 0.15cm\guillemotleft}%
      { \end{quotation}\vspace*{0.1cm}}



\newcommand{\subsubseccion}[1]{\subsubsection*{#1}\addcontentsline{toc}{subsubsection}{#1}}



% Circulos chicos: un cuarto, mitad, trescuartos y completo.
% \newcommand{\quarter}{
%                     \begin{tikzpicture}
%                     \fill[black] (0,0.11) arc (90:-90:0.1cm) ;
%                     \filldraw[white] [rotate=-90] (0,0.11) arc (90:-90:0.1cm);
%                     \draw (0,0) circle (0.1cm);
%                     \end{tikzpicture}
%                     }
% 
% \newcommand{\half}{
%                     \begin{tikzpicture}
%                     \draw (0,0) circle (0.1cm);
%                     \fill[black] (0,0.11) arc (90:-90:0.1cm) ;
%                     \end{tikzpicture}
%                    }
% 
% \newcommand{\threequarter}{
%                     \begin{tikzpicture}
%                     \fill[black] (0,0.10) arc (90:-90:0.1cm);
%                     \filldraw[black] [rotate=-90] (0,0.10) arc (90:-90:0.1cm);
%                     \draw (0,0) circle (0.1cm);
%                     \end{tikzpicture}
%                    }
% 
% 
% \newcommand{\full}{
%                     \begin{tikzpicture}
%                     \fill[black] (0,0.10) arc (180:-180:0.1cm);
%                     \end{tikzpicture}
%                      }



% Circulos grandes
\newcommand{\quarter}{
                    \begin{tikzpicture}
                    \fill[black] (0,0.16) arc (90:-90:0.9ex);
                    \filldraw[white] [rotate=-90] (0,0.16) arc (90:-90:0.9ex);
                    \draw (0,0) circle (1ex);
                    \end{tikzpicture}
                    }

\newcommand{\half}{
                    \begin{tikzpicture}
                    \fill[black] (0,0.165) arc (90:-90:0.9ex) ;
                    \draw (0,0) circle (1ex);
                    \end{tikzpicture}
                   }

\newcommand{\threequarter}{
                    \begin{tikzpicture}
                    \fill[black] (0,0.165) arc (90:-90:0.9ex) ;
                    \filldraw[black] [rotate=-90] (0,0.16) arc (90:-90:0.9ex) ;
                    \draw (0,0) circle (1ex);
                    \end{tikzpicture}
                        }


\newcommand{\full}{
                    \begin{tikzpicture}
                    \filldraw[black] [rotate=-90] (0,0.18) arc (180:-180:1ex) ;
                    \end{tikzpicture}
                     }
